% FILE: 00_project_identity.tex
% Project: schemas
% Thesis Role: data-schema-definitions-for-ma-and-receipt-surfaces

\section{Project Identity}
\label{sec:schemas-identity}

\subsection{Purpose}
\label{sec:schemas-identity-purpose}

The \texttt{schemas} project establishes the structural contract between two critical surfaces of the
process intelligence manufacturing pipeline: M\&A deck generation and cryptographic receipt
validation. Its primary artifact, \texttt{slide\_verification\_ledger.schema.json}, is a JSON Schema
draft-07 document that defines the data shape every cryptographic receipt must carry to be
admissible as board-level evidence. A companion schema, \texttt{otel\_trace\_integration\_schema.json},
governs how W3C OpenTelemetry trace telemetry is linked to BLAKE3 event-chain roots.

Without these schemas, the receipt chain is informally specified and cannot be mechanically
validated. With them, any receipt file can be tested for structural conformance before it is
admitted into the slide-to-receipt ledger. The schemas therefore occupy the lowest stratum of
the evidence stack: they precede all upstream claim taxonomy, rendering pipeline, and
cryptographic signing work.

The project is minimal by design. It contains exactly one schema in the canonical \texttt{schemas/}
directory, with no build system, no test runner attached at commit time, and no README. This
minimalism is not incidental---it reflects the deliberate choice to keep the contract surface
narrow and stable. Every field in the schema has a corresponding required entry in the five
production receipts that serve as the project's principal validation evidence.

\subsection{Repository Surface}
\label{sec:schemas-identity-surface}

The repository surface consists of two JSON Schema files occupying different directories:

\begin{itemize}
  \item \texttt{schemas/slide\_verification\_ledger.schema.json} --- the primary schema governing
    the \texttt{SlideVerificationLedger} type. This file uses \texttt{patternProperties} keyed on
    a UUIDv4 regular expression (\verb|^[0-9a-fA-F]{8}-...-[0-9a-fA-F]{12}$|) so that each
    key in the ledger object is itself a UUID identifying a specific M\&A deck slide.
  \item \texttt{standards/otel\_trace\_integration\_schema.json} --- a companion schema defining
    the \texttt{OTelTraceBlake3IntegrationSchema}. This schema sits in \texttt{standards/} rather
    than \texttt{schemas/}, a governance boundary that is currently undocumented and treated as
    an open question in the research program.
\end{itemize}

Both files reference JSON Schema draft-07 as their meta-schema. The \texttt{additionalProperties: false}
constraint on the ledger schema enforces strict conformance: any receipt carrying fields not listed
in the required set is structurally invalid and must be refused.

\subsection{Primary Research Role}
\label{sec:schemas-identity-role}

Within the dissertation, the \texttt{schemas} project plays the role of
\emph{data-schema-definitions-for-ma-and-receipt-surfaces}. It answers the question: what is the
minimal structural contract that a cryptographic receipt must satisfy to be admissible as
board-level process intelligence evidence?

The schemas do not perform computation. They do not discover process models, replay event logs, or
compute fitness metrics. Their function is entirely definitional: they assert what fields must
exist, what patterns those fields must match, and what enumerated values are permitted. In this
sense they are analogous to type signatures in a programming language---the contract that upstream
manufacturers and downstream consumers must both respect.

This definitional role places the schemas at the intersection of three dissertation themes:
(1)~the receipt-chain architecture that makes process claims cryptographically auditable;
(2)~the M\&A evidence stack that renders those claims board-admissible; and
(3)~the manufacturing discipline that refuses to ship artifacts that cannot pass a proof gate.

\subsection{Key Artifacts}
\label{sec:schemas-identity-artifacts}

\begin{description}
  \item[\texttt{slide\_verification\_ledger.schema.json}] The canonical schema. Eight required
    fields per ledger entry: \texttt{slide\_id}, \texttt{slide\_title}, \texttt{assertion\_text},
    \texttt{target\_log\_hash} (64-hex BLAKE3), \texttt{process\_model\_hash} (64-hex BLAKE3),
    \texttt{query\_definition}, \texttt{verification\_results} (with status enum
    \texttt{verified}~$|$~\texttt{failed}~$|$~\texttt{unverified}), and \texttt{validator\_signature}
    (128-hex Ed25519).
  \item[\texttt{otel\_trace\_integration\_schema.json}] Companion schema. Requires \texttt{trace\_id}
    (32-hex W3C TraceID), \texttt{event\_chain\_root} (64-hex BLAKE3), and a \texttt{spans} array
    where each span carries a \texttt{blake3\_receipt} computed as
    $\operatorname{BLAKE3}(\text{PriorReceipt} \| \text{SpanData})$.
  \item[\texttt{rec\_ebitda\_rework\_001.json} through \texttt{rec\_residual\_standard\_005.json}]
    Five production receipts in \texttt{receipts/} that conform to the ledger schema, carrying
    real UUIDv4 \texttt{slide\_id} values, 64-hex BLAKE3 hashes, wasm4pm query URIs,
    fitness/precision metrics in the 0.975--1.0 range, and 128-hex Ed25519 signatures prefixed
    \texttt{ed25519\_sig\_}.
\end{description}
